% ============================================================================
% 10-MINUTE BEAMER PRESENTATION
% ============================================================================
% Main-talk image files (place beside this .tex file, in figures/, or in
% Batch_Figures/):
%   nc_groupoverview.png
%   mci_groupoverview.png
%   nc_groupmisomaptocbv.png
%   mci_groupmisomaptocbv.png
%   nc_groupmisoco2tocbv.png
%   mci_groupmisoco2tocbv.png
%   statistical_frequency_model_nc_gain.png
%   statistical_frequency_model_mci_gain.png
%   statistical_model_comparison.png
%   statistical_group_comparison.png
%
% Optional backup images:
%   statistical_frequency_group_gain.png
%   statistical_frequency_model_nc_phase.png
%   statistical_frequency_model_mci_phase.png
%   statistical_frequency_group_coherence.png
%   statistical_frequency_group_phase.png
%   statistical_input_association.png
%   statistical_pathway_balance.png
%
% The \resultimage command checks beside this file, figures/, and
% Batch_Figures/. If an image is missing, the compiled slide shows a visible
% placeholder with the required filename.
%
% Timing: approximately 14--15 minutes for the 18 main frames, including the
% title. Everything after \appendix is optional backup.
% ============================================================================

\documentclass[aspectratio=169,11pt]{beamer}

\usepackage[T1]{fontenc}
\usepackage[utf8]{inputenc}
\usepackage{amsmath,amssymb}
\usepackage{booktabs}
\usepackage{array}
\usepackage{graphicx}
\usepackage{xcolor}
\usepackage{hyperref}

\graphicspath{{./}{figures/}{Batch_Figures/}}

% ---------------------------------------------------------------------------
% Visual style
% ---------------------------------------------------------------------------

\definecolor{ProjectBlue}{RGB}{0,114,189}
\definecolor{ProjectOrange}{RGB}{217,83,25}
\definecolor{ProjectNavy}{RGB}{25,55,85}
\definecolor{ProjectLightBlue}{RGB}{232,243,250}
\definecolor{ProjectLightOrange}{RGB}{252,238,228}
\definecolor{ProjectGray}{RGB}{90,95,100}
\definecolor{ProjectLightGray}{RGB}{244,245,246}

\setbeamercolor{normal text}{fg=black,bg=white}
\setbeamercolor{structure}{fg=ProjectBlue}
\setbeamercolor{frametitle}{fg=white,bg=ProjectNavy}
\setbeamercolor{title}{fg=ProjectNavy}
\setbeamercolor{subtitle}{fg=ProjectGray}
\setbeamercolor{block title}{fg=white,bg=ProjectBlue}
\setbeamercolor{block body}{fg=black,bg=ProjectLightBlue}
\setbeamercolor{block title alerted}{fg=white,bg=ProjectOrange}
\setbeamercolor{block body alerted}{fg=black,bg=ProjectLightOrange}
\setbeamercolor{footline}{fg=ProjectGray,bg=white}

\setbeamertemplate{navigation symbols}{}
\setbeamertemplate{itemize item}{\color{ProjectBlue}$\blacktriangleright$}
\setbeamertemplate{itemize subitem}{\color{ProjectOrange}$\bullet$}
\setbeamertemplate{enumerate item}{\color{ProjectBlue}\insertenumlabel.}
\setbeamertemplate{frametitle}{%
    \nointerlineskip
    \begin{beamercolorbox}[wd=\paperwidth,ht=0.78cm,dp=0.22cm,leftskip=0.42cm]{frametitle}
        \usebeamerfont{frametitle}\insertframetitle
    \end{beamercolorbox}}
\setbeamertemplate{footline}{%
    \leavevmode
    \hbox{%
        \begin{beamercolorbox}[wd=0.84\paperwidth,ht=2.5ex,dp=1.2ex,leftskip=0.4cm]{footline}
            MISO transfer function analysis of cerebral hemodynamics
        \end{beamercolorbox}%
        \begin{beamercolorbox}[wd=0.16\paperwidth,ht=2.5ex,dp=1.2ex,rightskip=0.4cm plus1fil]{footline}
            \hfill\insertframenumber
        \end{beamercolorbox}}
    \vskip0pt}

\setlength{\leftmargini}{0.24in}
\setlength{\leftmarginii}{0.20in}

% ---------------------------------------------------------------------------
% Reusable commands
% ---------------------------------------------------------------------------

\newcommand{\PETCOtwo}{P_{\mathrm{ET}}CO_2}
\newcommand{\COtwo}{CO_2}

% Search beside the .tex file, in figures/, and in Batch_Figures/.
\newcommand{\resultimage}[2]{%
    \IfFileExists{#1}{%
        \includegraphics[width=\linewidth,height=#2,keepaspectratio]{#1}%
    }{%
        \IfFileExists{figures/#1}{%
            \includegraphics[width=\linewidth,height=#2,keepaspectratio]{figures/#1}%
        }{%
            \IfFileExists{Batch_Figures/#1}{%
                \includegraphics[width=\linewidth,height=#2,keepaspectratio]{Batch_Figures/#1}%
            }{%
                \fcolorbox{ProjectGray}{ProjectLightGray}{%
                    \parbox[c][#2][c]{0.92\linewidth}{%
                        \centering
                        \Large\color{ProjectGray}\textbf{FIGURE PLACEHOLDER}\par\medskip
                        \normalsize Upload this file beside the presentation,
                        into \texttt{figures/}, or into
                        \texttt{Batch\_Figures/}:\par\medskip
                        \texttt{\detokenize{#1}}}}%
            }%
        }%
    }%
}

\newcommand{\takeaway}[1]{%
    \begin{beamercolorbox}[rounded=true,sep=0.12cm,wd=\linewidth]{block body alerted}
        \centering\textbf{Takeaway:} #1
    \end{beamercolorbox}}

\newcommand{\methodbox}[2]{%
    \fcolorbox{ProjectBlue}{ProjectLightBlue}{%
        \parbox[c][1.12cm][c]{0.155\linewidth}{%
            \centering\textbf{#1}\par\scriptsize #2}}}

% ---------------------------------------------------------------------------
% Title information
% ---------------------------------------------------------------------------

\title{Separating MAP- and CO$_2$-Associated Contributions to Cerebral Blood
Flow Velocity}
\subtitle{A multivariable transfer function analysis in mild cognitive impairment}
\author{Amogh Naganand}
\date{July 21, 2026}

\begin{document}

% ============================================================================
% MAIN TALK: 18 FRAMES INCLUDING TITLE / APPROXIMATELY 14--15 MINUTES
% ============================================================================

% TALK TRACK (0:00--0:30)
% Introduce the physiological question in one sentence.
\begin{frame}[plain]
    \titlepage

    \vspace{-0.3cm}
    \begin{center}
        \color{ProjectNavy}
        \textbf{Core question:} Does accounting for spontaneous CO$_2$
        change how we interpret MAP--CBFV dynamics in NC and MCI?
    \end{center}
\end{frame}

% TALK TRACK (0:30--1:20)
% Explain that CBFV reflects multiple overlapping influences. Avoid causal language.
\begin{frame}{Why this question matters}
    \begin{columns}[T,onlytextwidth]
        \begin{column}{0.31\textwidth}
            \begin{block}{Pressure}
                MAP fluctuations are the standard input used to assess
                dynamic cerebral autoregulation.
            \end{block}
        \end{column}
        \begin{column}{0.31\textwidth}
            \begin{alertblock}{Carbon dioxide}
                CO$_2$ is vasoactive and changes cerebral vascular tone in
                overlapping frequency ranges.
            \end{alertblock}
        \end{column}
        \begin{column}{0.31\textwidth}
            \begin{block}{Cerebral output}
                Transcranial Doppler CBFV reflects the combined behavior of
                pressure, CO$_2$, and unmeasured mechanisms.
            \end{block}
        \end{column}
    \end{columns}

    \vspace{0.35cm}
    \begin{center}
        \Large
        \color{ProjectBlue}\textbf{MAP}
        \color{black}$\quad+\quad$
        \color{ProjectOrange}\textbf{CO$_2$}
        \color{black}$\quad\longrightarrow\quad$
        \color{ProjectNavy}\textbf{CBFV}
    \end{center}

    \vspace{0.25cm}
    \takeaway{A MAP-only model may mix pressure-associated dynamics with
    variation shared with CO$_2$.}
\end{frame}

% TALK TRACK (1:20--2:30)
% Define every subscript directly so the audience does not need abstract x1, x2, or y notation.
\begin{frame}{SISO and MISO answer different questions}
    \begin{columns}[T,onlytextwidth]
        \begin{column}{0.43\textwidth}
            \begin{block}{Separate SISO models}
                \centering
                \[
                    H_{\mathrm{MAP}}^{\mathrm{SISO}}
                    =\frac{S_{\mathrm{MAP},\mathrm{CBFV}}}
                    {S_{\mathrm{MAP},\mathrm{MAP}}}
                \]
                \[
                    H_{\COtwo}^{\mathrm{SISO}}
                    =\frac{S_{\COtwo,\mathrm{CBFV}}}
                    {S_{\COtwo,\COtwo}}
                \]
                \small Pairwise input--output relationships
            \end{block}
        \end{column}

        \begin{column}{0.53\textwidth}
            \begin{alertblock}{One simultaneous MISO model}
                \centering
                \resizebox{0.96\linewidth}{!}{$
                \begin{bmatrix}
                    S_{\mathrm{MAP},\mathrm{MAP}} &
                    S_{\mathrm{MAP},\COtwo}\\
                    S_{\COtwo,\mathrm{MAP}} &
                    S_{\COtwo,\COtwo}
                \end{bmatrix}
                \begin{bmatrix}
                    H_{\mathrm{MAP}}^{\mathrm{MISO}}\\
                    H_{\COtwo}^{\mathrm{MISO}}
                \end{bmatrix}
                =
                \begin{bmatrix}
                    S_{\mathrm{MAP},\mathrm{CBFV}}\\
                    S_{\COtwo,\mathrm{CBFV}}
                \end{bmatrix}
                $}

                \medskip
                \small Each pathway is estimated while accounting for the
                other measured input
            \end{alertblock}
        \end{column}
    \end{columns}

    \vspace{0.2cm}
    \textbf{Central hypothesis:} Simultaneously modeling MAP and CO$_2$ will
    produce different pathway-specific transfer function estimates than
    modeling each input separately. These conditional estimates may also
    reveal NC--MCI differences.
\end{frame}

% TALK TRACK (2:30--3:05)
% Keep the workflow on one line. The scientific results receive more time than software details.
\begin{frame}{Pipeline}
    \centering
    \resizebox{0.98\textwidth}{!}{%
        \methodbox{Input}{single, batch, or synthetic}
        \hspace{0.08cm}$\longrightarrow$\hspace{0.08cm}
        \methodbox{Preprocess}{clean, resample, normalize}
        \hspace{0.08cm}$\longrightarrow$\hspace{0.08cm}
        \methodbox{Welch spectra}{one shared spectral estimate}
        \hspace{0.08cm}$\longrightarrow$\hspace{0.08cm}
        \methodbox{Models}{SISO and MISO}
        \hspace{0.08cm}$\longrightarrow$\hspace{0.08cm}
        \methodbox{Output}{plots, statistics, Excel}}

    \vspace{0.65cm}
    \begin{block}{Consistent analysis path}
        The same preprocessing and Welch spectra are passed into both models.
        Subject results are retained for group means, variability, paired model
        comparisons, and independent NC--MCI comparisons.
    \end{block}

    \vspace{0.25cm}
    \takeaway{Model differences reflect the transfer-function formulation,
    rather than separate preprocessing or spectral estimates.}
\end{frame}

% TALK TRACK (3:05--3:55)
% State only the settings that influence interpretation. Do not list every toggle.
\begin{frame}{Dataset and analysis settings}
    \begin{columns}[T,onlytextwidth]
        \begin{column}{0.49\textwidth}
            \begin{block}{Current cohort}
                \begin{itemize}
                    \item \textcolor{ProjectBlue}{\textbf{23 NC}} participants
                    \item \textcolor{ProjectOrange}{\textbf{52 MCI}} participants
                    \item All 75 completed the current TFA workflow
                    \item Similar age, sex, MAP, $\PETCOtwo$, CBFV, and
                    recording duration
                    \item MoCA was lower in MCI, as expected
                \end{itemize}
            \end{block}
        \end{column}

        \begin{column}{0.49\textwidth}
            \begin{alertblock}{Core settings}
                \begin{itemize}
                    \item Resampled to 4 Hz
                    \item CBFV normalized to percent baseline
                    \item 128-s Hann windows, 50\% overlap
                    \item Minimum of 3 Welch windows
                    \item Frequency range: 0--0.35 Hz
                    \item Bands: VVLF, VLF, LF, HF
                    \item Same smoothing kernel for every spectrum
                \end{itemize}
            \end{alertblock}
        \end{column}
    \end{columns}

    \takeaway{With these settings, three Welch windows require at least
    1,024 samples, or 256 seconds at 4 Hz.}
\end{frame}

% TALK TRACK (3:55--4:30)
% Establish the spectra and input relationship before presenting transfer functions.
\begin{frame}{NC signal and spectral overview}
    \centering
    \resultimage{nc_groupoverview.png}{0.82\textheight}
\end{frame}

% TALK TRACK (4:30--5:05)
% Use the same four panels to orient the audience to the MCI input and output data.
\begin{frame}{MCI signal and spectral overview}
    \centering
    \resultimage{mci_groupoverview.png}{0.82\textheight}
\end{frame}

% TALK TRACK (5:05--5:55)
% Define what the audience will see in the next two result slides.
\begin{frame}{Reading a transfer-function result}
    \begin{columns}[T,onlytextwidth]
        \begin{column}{0.31\textwidth}
            \begin{block}{Gain}
                \centering
                \Large $\left|H(f)\right|$

                \normalsize\medskip
                Magnitude of the frequency-specific CBFV response associated
                with an input oscillation.
            \end{block}
        \end{column}
        \begin{column}{0.31\textwidth}
            \begin{block}{Phase}
                \centering
                \Large $\arg\!\left[H(f)\right]$

                \normalsize\medskip
                Timing relationship between the input and CBFV. Group phase
                uses circular means and circular SD.
            \end{block}
        \end{column}
        \begin{column}{0.31\textwidth}
            \begin{alertblock}{Coherence}
                Strength of the linear frequency-specific relationship.

                \medskip
                SISO uses ordinary coherence; MISO reports pathway-specific
                partial coherence and combined multiple coherence.
            \end{alertblock}
        \end{column}
    \end{columns}

    \vspace{0.35cm}
    \takeaway{The complete transfer function is frequency-dependent; band
    averages summarize these curves but do not replace them.}
\end{frame}

% TALK TRACK (5:55--6:35)
% Full-size primary result for the NC MAP pathway.
\begin{frame}{NC: MISO MAP-to-CBFV transfer function}
    \centering
    \resultimage{nc_groupmisomaptocbv.png}{0.82\textheight}
\end{frame}

% TALK TRACK (6:35--7:15)
% Full-size primary result for the MCI MAP pathway.
\begin{frame}{MCI: MISO MAP-to-CBFV transfer function}
    \centering
    \resultimage{mci_groupmisomaptocbv.png}{0.82\textheight}
\end{frame}

% TALK TRACK (7:15--7:55)
% Full-size primary result for the NC CO2 pathway.
\begin{frame}{NC: MISO CO$_2$-to-CBFV transfer function}
    \centering
    \resultimage{nc_groupmisoco2tocbv.png}{0.82\textheight}
\end{frame}

% TALK TRACK (7:55--8:35)
% Full-size primary result for the MCI CO2 pathway.
\begin{frame}{MCI: MISO CO$_2$-to-CBFV transfer function}
    \centering
    \resultimage{mci_groupmisoco2tocbv.png}{0.82\textheight}
\end{frame}

% TALK TRACK (8:35--9:25)
% Now that the audience has seen the transfer functions, explain how comparisons are tested.
\begin{frame}{Comparing the Models}
    \small
    \begin{columns}[T,onlytextwidth]
        \begin{column}{0.48\textwidth}
            \begin{block}{MISO versus SISO}
                For each participant and result:

                \[
                    d=\text{MISO value}-\text{SISO value}
                \]

                \[
                    H_0:\ \text{average difference}=0.
                \]
            \end{block}
        \end{column}
        \begin{column}{0.48\textwidth}
            \begin{block}{NC versus MCI}
                For each MISO result, participant values in NC are compared
                with participant values in MCI.

                \[
                    H_0:\ \overline{x}_{\mathrm{NC}}
                    =\overline{x}_{\mathrm{MCI}}.
                \]
            \end{block}
        \end{column}
    \end{columns}

    \vspace{0.20cm}
    \begin{alertblock}{Why adjust the P values?}
        Testing many results increases the chance of a false-positive finding.
        Benjamini--Hochberg adjustment accounts for these multiple tests;
        adjusted $P<0.05$ is treated as evidence of a difference.
    \end{alertblock}
\end{frame}

% TALK TRACK (9:25--10:20)
% Read each column from top to bottom: curves, paired difference, adjusted P value.
\begin{frame}{NC: frequency-wise MISO versus SISO gain}
    \centering
    \resultimage{statistical_frequency_model_nc_gain.png}{0.73\textheight}

    \vspace{0.08cm}
    \takeaway{NC MAP gain changed relatively little, whereas CO$_2$ gain
    differed across contiguous mid- and higher-frequency regions.}
\end{frame}

% TALK TRACK (10:20--11:15)
% Emphasize the broad CO2 model effect without calling either model ground truth.
\begin{frame}{MCI: frequency-wise MISO versus SISO gain}
    \centering
    \resultimage{statistical_frequency_model_mci_gain.png}{0.73\textheight}

    \vspace{0.08cm}
    \takeaway{The MCI CO$_2$ pathway showed a broad model effect; MAP gain
    differences were smaller and concentrated mainly at the lowest frequencies.}
\end{frame}

% TALK TRACK (11:15--12:10)
% This is the primary subject-level inference. Point to paired direction and adjusted P values.
\begin{frame}{Primary band-level evidence for the model effect}
    \centering
    \resultimage{statistical_model_comparison.png}{0.69\textheight}

    \vspace{0.08cm}
    \takeaway{MISO MAP gain was lower in VVLF in both groups, while the
    larger and more consistent model-related changes occurred in CO$_2$ gain.}
\end{frame}

% TALK TRACK (12:10--13:05)
% The primary biological comparison is MISO NC versus MCI at the band level.
\begin{frame}{NC versus MCI: band-level comparisons}
    \centering
    \resultimage{statistical_group_comparison.png}{0.69\textheight}

    \vspace{0.08cm}
    \takeaway{None of the tested NC--MCI band differences had an adjusted
    $P<0.05$. The current data therefore did not identify a robust group
    difference in these MISO metrics.}
\end{frame}

% TALK TRACK (13:05--14:10)
% End with what is supported, what is not established, and the next validation steps.
\begin{frame}{Conclusions and next steps}
    \begin{enumerate}
        \item \textbf{Model comparison:} MISO and SISO produced statistically
        different pathway-specific estimates in several comparisons,
        particularly for CO$_2$ gain.

        \item \textbf{Group comparison:} The tested MISO metrics did not
        identify a robust difference between NC and MCI in the current data.

        \item \textbf{Interpretation:} The results support considering
        multiple measured inputs when estimating autoregulatory dynamics, but
        adding CO$_2$ alone was not sufficient to establish an NC--MCI marker.
        Additional inputs, smaller effects, and biological variability remain
        possible explanations requiring further investigation.
    \end{enumerate}

    \vspace{0.25cm}
    \begin{alertblock}{Next steps}
        \begin{itemize}
            \item Obtain longer recordings that permit meaningful
            cross-validation of SISO and MISO predictions.
            \item Complete the manuscript and prepare it for journal
            submission or preprint release.
        \end{itemize}
    \end{alertblock}

\end{frame}

% ============================================================================
% BACKUP SLIDES -- NOT PART OF THE 10-MINUTE TALK
% ============================================================================

\appendix

\begin{frame}{Backup: frequency-band definitions}
    \centering
    \begin{tabular}{lll}
        \toprule
        Band & Frequency interval & Included frequency bins \\
        \midrule
        VVLF & $0.005 \leq f < 0.024$ Hz & 0.0078--0.0234 Hz \\
        VLF  & $0.024 \leq f < 0.070$ Hz & 0.0313--0.0625 Hz \\
        LF   & $0.070 \leq f < 0.200$ Hz & 0.0703--0.1953 Hz \\
        HF   & $0.200 \leq f \leq 0.350$ Hz & 0.2031--0.3438 Hz \\
        \bottomrule
    \end{tabular}

    \vspace{0.55cm}
    Frequency resolution:
    \[
        \Delta f=\frac{4\ \mathrm{Hz}}{512}=0.0078125\ \mathrm{Hz}.
    \]

    Band averages were calculated within subject first, so each subject
    contributed one value per band to the formal band-level tests.
\end{frame}

\begin{frame}{Backup: SISO and MISO derivation}
    \small
    For MAP as one input,
    \[
        \mathrm{CBFV}=H_{\mathrm{MAP}}\,\mathrm{MAP}+E,
        \qquad
        S_{\mathrm{MAP},\mathrm{CBFV}}
        =H_{\mathrm{MAP}}S_{\mathrm{MAP},\mathrm{MAP}},
        \qquad
        H_{\mathrm{MAP}}
        =\frac{S_{\mathrm{MAP},\mathrm{CBFV}}}
        {S_{\mathrm{MAP},\mathrm{MAP}}}.
    \]

    With MAP and CO$_2$ entered simultaneously,
    \[
        \mathrm{CBFV}
        =H_{\mathrm{MAP}}\,\mathrm{MAP}
        +H_{\COtwo}\,\COtwo+E,
    \]
    \resizebox{0.88\linewidth}{!}{$
        \underbrace{\begin{bmatrix}
            S_{\mathrm{MAP},\mathrm{MAP}} & S_{\mathrm{MAP},\COtwo}\\
            S_{\COtwo,\mathrm{MAP}} & S_{\COtwo,\COtwo}
        \end{bmatrix}}_{\mathbf S_{\mathrm{inputs}}}
        \underbrace{\begin{bmatrix}
            H_{\mathrm{MAP}}^{\mathrm{MISO}}\\
            H_{\COtwo}^{\mathrm{MISO}}
        \end{bmatrix}}_{\mathbf H}
        =
        \underbrace{\begin{bmatrix}
            S_{\mathrm{MAP},\mathrm{CBFV}}\\
            S_{\COtwo,\mathrm{CBFV}}
        \end{bmatrix}}_{\mathbf S_{\mathrm{input},\mathrm{CBFV}}}
    $}

    \medskip
    The implementation solves this spectral system at each frequency using
    MATLAB matrix left division. The current analysis is unregularized.
\end{frame}

\begin{frame}{Backup: how phase was summarized}
    \begin{columns}[T,onlytextwidth]
        \begin{column}{0.47\textwidth}
            \begin{block}{Circular mean}
                \[
                    C=\frac{1}{n}\sum_i\cos\theta_i,
                    \qquad
                    S=\frac{1}{n}\sum_i\sin\theta_i
                \]
                \[
                    \overline\theta=\operatorname{atan2}(S,C)
                \]
            \end{block}
        \end{column}
        \begin{column}{0.47\textwidth}
            \begin{alertblock}{Circular spread}
                \[
                    R=\sqrt{C^2+S^2}
                \]
                \[
                    s_{\mathrm{circ}}=\sqrt{-2\ln R}
                \]
            \end{alertblock}
        \end{column}
    \end{columns}

    \vspace{0.35cm}
    Statistics used wrapped phase. For display, the circular group mean was
    unwrapped across frequency and retained the same circular SD.
\end{frame}

\begin{frame}{Backup: what a BH-adjusted P value means}
    \begin{block}{One raw P value}
        Tests one null hypothesis without accounting for the other tests in
        the family.
    \end{block}

    \begin{alertblock}{BH-adjusted P value}
        Accounts for the number and ranking of related P values. If
        $P_{\mathrm{BH}}<0.05$, that test is significant while controlling the
        false discovery rate within the defined family.
    \end{alertblock}

    \begin{itemize}
        \item $P_{\mathrm{BH}}>0.05$ does \textbf{not} prove that two curves or
        groups are the same.
        \item More subjects can improve power if the underlying effect is
        consistent, but sample size does not guarantee significance.
        \item Standard deviation affects P values because it affects the
        uncertainty of the estimated difference.
    \end{itemize}
\end{frame}

\begin{frame}{Backup: NC frequency-wise phase comparison}
    \centering
    \resultimage{statistical_frequency_model_nc_phase.png}{0.75\textheight}

    \takeaway{Phase comparisons use circular paired differences and
    permutation tests; the displayed group means are circular means.}
\end{frame}

\begin{frame}{Backup: MCI frequency-wise phase comparison}
    \centering
    \resultimage{statistical_frequency_model_mci_phase.png}{0.75\textheight}

    \takeaway{The same circular procedure is used in MCI, with BH adjustment
    performed within each frequency-comparison curve.}
\end{frame}

\begin{frame}{Backup: NC--MCI frequency-wise gain}
    \centering
    \resultimage{statistical_frequency_group_gain.png}{0.75\textheight}

    \takeaway{The isolated adjusted result at 0.1484 Hz is exploratory and
    was not accompanied by a robust prespecified band-level group difference.}
\end{frame}

\begin{frame}{Backup: NC--MCI frequency-wise coherence}
    \centering
    \resultimage{statistical_frequency_group_coherence.png}{0.75\textheight}

    \takeaway{Frequency-wise group coherence comparisons are exploratory;
    primary inference uses subject-level band values.}
\end{frame}

\begin{frame}{Backup: NC--MCI frequency-wise phase}
    \centering
    \resultimage{statistical_frequency_group_phase.png}{0.75\textheight}

    \takeaway{A high adjusted P value means the available data do not provide
    evidence of a group difference at that frequency.}
\end{frame}

\begin{frame}{Backup: input coherence and gain change}
    \centering
    \resultimage{statistical_input_association.png}{0.73\textheight}

    \takeaway{Several low-frequency associations were significant, but they
    are exploratory because both variables come from the same spectral matrix.}
\end{frame}

\begin{frame}{Backup: MAP--CO$_2$ pathway balance}
    \centering
    \resultimage{statistical_pathway_balance.png}{0.73\textheight}

    \takeaway{MAP partial coherence became larger than CO$_2$ partial
    coherence above VVLF, with no corrected NC--MCI difference in balance.}
\end{frame}

\begin{frame}{Backup: numerical conditioning}
    \begin{columns}[T,onlytextwidth]
        \begin{column}{0.48\textwidth}
            \begin{block}{Current observations}
                \begin{itemize}
                    \item Median HF condition number: 10.9 NC, 15.4 MCI
                    \item HF values $>100$: 60/437 NC, 126/988 MCI
                    \item Maximum: 7,467 NC, 17,366 MCI
                    \item High values often occurred with low CO$_2$ power or
                    large input-power imbalance
                \end{itemize}
            \end{block}
        \end{column}
        \begin{column}{0.48\textwidth}
            \begin{alertblock}{Interpretation}
                \begin{itemize}
                    \item The raw condition number is scale-dependent.
                    \item A high value does not automatically prove
                    physiological collinearity.
                    \item It does indicate that some estimates may be
                    sensitive to small spectral changes.
                    \item Ridge regularization remains a planned sensitivity
                    decision, not part of the reported solve.
                \end{itemize}
            \end{alertblock}
        \end{column}
    \end{columns}
\end{frame}

\begin{frame}{Backup: key numeric model comparisons}
    \scriptsize
    \centering
    \begin{tabular}{lllrrrr}
        \toprule
        Group & Path & Band & MISO & SISO & Difference & BH $P$ \\
        \midrule
        NC  & MAP & VVLF & 1.092 & 1.332 & $-0.240$ & 0.0052 \\
        MCI & MAP & VVLF & 0.837 & 1.020 & $-0.183$ & 0.0038 \\
        NC  & MAP & LF   & 2.111 & 2.046 & $+0.065$ & 0.0464 \\
        NC  & CO$_2$ & LF & 1.761 & 2.498 & $-0.737$ & 0.00052 \\
        NC  & CO$_2$ & HF & 2.695 & 4.944 & $-2.250$ & 0.00017 \\
        MCI & CO$_2$ & VLF & 2.066 & 2.850 & $-0.784$ & 0.000083 \\
        MCI & CO$_2$ & LF & 1.599 & 2.892 & $-1.293$ & 0.0000021 \\
        MCI & CO$_2$ & HF & 2.823 & 5.458 & $-2.635$ & 0.0052 \\
        \bottomrule
    \end{tabular}

    \vspace{0.45cm}
    Gain units are percentage points of baseline CBFV per mmHg. Differences
    are MISO minus SISO. BH adjustment covered the 16 primary gain tests.
\end{frame}

\end{document}
